% !TEX program = latexmk (xelatex)
\documentclass[12pt, a4paper]{article}

\usepackage[polish]{babel}
\usepackage[utf8]{inputenc}
\usepackage[T1]{fontenc}
\usepackage{lato}
\renewcommand{\familydefault}{\sfdefault}
\setlength{\parskip}{0.3em}
\usepackage{fancyhdr}
\usepackage{tikz}
\usetikzlibrary{arrows.meta}
\usepackage{hyperref}
\usepackage{listings}
\usepackage{xcolor}
\usepackage{float}

% Konfiguracja wyświetlania kodu i logów
\lstset{
    basicstyle=\ttfamily\small,
    breaklines=true,
    frame=single,
    backgroundcolor=\color{gray!10},
    keywordstyle=\color{blue},
    commentstyle=\color{green!50!black},
    stringstyle=\color{red},
    showstringspaces=false,
    extendedchars=true,
    literate={ą}{{\k{a}}}1 {ć}{{\'c}}1 {ę}{{\k{e}}}1 {ł}{{\l{}}}1 {ń}{{\'n}}1 {ó}{{\'o}}1 {ś}{{\'s}}1 {ź}{{\'z}}1 {ż}{{\.z}}1
             {Ą}{{\k{A}}}1 {Ć}{{\'C}}1 {Ę}{{\k{E}}}1 {Ł}{{\L{}}}1 {Ń}{{\'N}}1 {Ó}{{\'O}}1 {Ś}{{\'S}}1 {Ź}{{\'Z}}1 {Ż}{{\.Z}}1
}

\hypersetup{
    colorlinks=true,
    linkcolor=black,
    urlcolor=blue
}

\begin{document}
\pagestyle{myheadings}
\markright{PSI Dokumentacja końcowa projektu \hfill \today \hfill}

\noindent{\large{Zespół 76}}\\

\vspace{-0.5em}
\noindent Wojciech Kukiełka 331394\\
Tomasz Okoń 331414\\
Adam Szkolaski 325237

\hrulefill

\begin{center}
    \LARGE \textbf{System anonimizujący ,,Mini-TOR''}
\end{center}

\tableofcontents
\newpage

\section{Opis projektu i implementacji}

Celem projektu było stworzenie systemu anonimizującego ruch sieciowy TCP, inspirowanego rozwiązaniem Tor, lecz w uproszczonej architekturze opartej na pojedynczym węźle pośredniczącym (Proxy Node). System ma za zadanie ukryć tożsamość klienta przed serwerem docelowym oraz utrudnić analizę statystyczną ruchu poprzez mechanizmy opóźnień i segmentacji.

\subsection{Architektura systemu}
Komunikacja odbywa się w następującej architekturze:
\begin{enumerate}
    \item \textbf{Klient (Client):} Aplikacja korzystająca z naszej biblioteki pythonowej, która zestawia bezpieczny tunel TLS do węzła pośredniczącego.
    \item \textbf{Węzeł pośredniczący (Node):} Serwer przyjmujący połączenia od klientów, deszyfrujący pierwszą warstwę, a następnie nawiązujący połączenie z~serwerem docelowym. Węzeł odpowiada za anonimizację (podmianę adresu źródłowego) oraz kształtowanie ruchu (Traffic Shaping).
    \item \textbf{Serwer docelowy (Target Server):} Odbiorca ruchu, który widzi jedynie adres IP węzła pośredniczącego.
\end{enumerate}

\begin{figure}[h!]
\centering
\begin{tikzpicture}
    \path (0,0)  node[circle, draw, align=center] {Klient\\(Biblioteka)}
          (5, 0) node[circle, draw, align=center] {Węzeł\\(Node)}
          (10, 0) node[circle, draw, align=center] {Serwer\\Docelowy};
    \draw[arrows={- Latex[length=8pt]}] (1.2, 0) -- node[above] {Double TLS} node[below] {\scriptsize (CONNECT)} (3.8, 0);
    \draw[arrows={- Latex[length=8pt]}] (6.2, 0) -- node[above] {TLS} node[below] {\scriptsize (Traffic Shaping)} (8.6, 0);
\end{tikzpicture}
\caption{Architektura rozwiązania}
\label{fig:arch}
\end{figure}

\subsection{Zastosowane technologie}
Implementacja została wykonana w języku \textbf{Python 3.11}, z wykorzystaniem biblioteki standardowej \texttt{ssl} do obsługi szyfrowania oraz modułu \texttt{socket} do obsługi gniazd. Całość środowiska jest orkiestrowana przy pomocy \textbf{Docker Compose}. Do generowania certyfikatów X.509 wykorzystano pakiet \textbf{OpenSSL}.

\section{Scenariusze komunikacji}

\subsection{SC1: Nawiązanie połączenia i transfer danych}
\textbf{Aktorzy:} Klient, Węzeł, Serwer Docelowy.
\begin{enumerate}
    \item Klient inicjuje połączenie TLS do Węzła na ustalonym porcie.
    \item Klient wysyła wewnątrz tunelu TLS komunikat sterujący \texttt{CONNECT} wskazujący adres
          i port serwera docelowego.
    \item Węzeł parsuje komunikat i nawiązuje \textbf{zwykłe} połączenie TCP z serwerem docelowym.
    \item Po sukcesie, Węzeł przechodzi w tryb przekazywania danych (transparent proxy) z włączonym mechanizmem losowych opóźnień.
    \item Poprzez zestawiony tunel, klient inicjuje sesję TLS z serwerem a następnie z nim rozmawia.
\end{enumerate}

\subsection{SC2: Obsługa błędów połączenia}
\textbf{Sytuacja:} Klient żąda połączenia z nieosiągalnym serwerem.
\begin{enumerate}
    \item Klient wysyła żądanie \texttt{CONNECT} do nieistniejącego hosta.
    \item Węzeł podejmuje próbę połączenia, która kończy się niepowodzeniem (timeout lub connection refused).
    \item Węzeł wysyła komunikat ERROR.
    \item Biblioteka po stronie klienta zgłasza wyjątek \texttt{ConnectionError}.
\end{enumerate}

Co ciekawe, nie udało nam się zaobserwować tego scenariusza. Nawet gdy kontener serwera
nie był włączony, węzłowi udawało się nawiązać połączenie z nieistniejącą usługą. Błąd występował
dopiero po tym jak klient zaczął TLS handshake z serwerem. Tak działa docker -- gniazda akceptują
połączenia, nawet jak nie odbiera ich żaden kontener.

\section{Definicje komunikatów}

Wymiana danych sterujących odbywa się wewnątrz zestawionego tunelu TLS. Protokół jest tekstowy i uproszczony względem HTTP/SOCKS.

\subsection{Żądanie połączenia (Od Klienta do Węzła)}
Pierwszy komunikat wysyłany przez klienta po zestawieniu tunelu TLS.

\begin{lstlisting}[language=bash, caption=Format komunikatu CONNECT]
CONNECT <target_host>:<target_port>\n
\end{lstlisting}

\noindent \textbf{Pola:}
\begin{itemize}
    \item \texttt{CONNECT} – słowo kluczowe.
    \item \texttt{target\_host} – nazwa domenowa lub adres IP serwera docelowego.
    \item \texttt{target\_port} – port docelowy (np. 80).
    \item \texttt{\textbackslash n} – znak nowej linii kończący nagłówek.
\end{itemize}

\subsection{Strumień danych (Dwukierunkowy)}
Po poprawnym przetworzeniu nagłówka \texttt{CONNECT} wszelkie kolejne bajty są traktowane jako `payload` (dane użytkownika).
\begin{itemize}
    \item \textbf{Węzeł $\rightarrow$ Serwer:} Dane są wysyłane po segmentacji i wprowadzeniu losowych opóźnień (algorytm \textit{shaper}).
    \item \textbf{Serwer $\rightarrow$ Węzeł $\rightarrow$ Klient:} Odpowiedzi wracają tą samą drogą.
\end{itemize}

\section{Opisy zachowania podmiotów komunikacji}

\subsection{Klient (Biblioteka)}
Biblioteka udostępnia interfejs gniazda. Jej zadaniem jest ukrycie przed programistą faktu, że połączenie nie jest bezpośrednie. Odpowiada za wczytanie certyfikatu CA (do weryfikacji węzła), zestawienie tunelu SSL/TLS oraz wysłanie nagłówka protokołu.

\subsection{Węzeł Pośredniczący (Shaper & Proxy)}
Serce systemu. Działa w pętli zdarzeń, obsługując wiele połączeń jednocześnie.

\begin{minipage}{1.0\linewidth}
    Kluczowa logika \textit{Traffic Shaping}:
    \begin{itemize}
        \item Buforuje dane przychodzące od klienta.
        \item Dzieli bufor na losowe fragmenty (chunks).
        \item Wstrzymuje wysyłanie każdego fragmentu o losowy czas $t \in [0.05s, 0.2s]$.
        \item Uniemożliwia korelację czasową pakietów przychodzących i wychodzących.
    \end{itemize}
\end{minipage}

Aby te manipulacje powyżej funkcji \texttt{send} przekładały się na ruch sieciowy, wyłączyliśmy
algorytm Nagle'a.

\subsection{Serwer Docelowy}
Pasywny podmiot, który loguje przychodzące połączenia. W kontekście projektu służy do weryfikacji, czy adres źródłowy faktycznie należy do Węzła, a nie do Klienta.

\section{Wyniki testów}

Testy przeprowadzono w środowisku skonteneryzowanym (Docker Compose). Poniżej przedstawiono dowody poprawnego działania systemu.

\subsection{Test funkcjonalny: Przesył danych}
Przykład poprawnego wykonania demonstracji protokołu.

\begin{lstlisting}[caption=Logi z klienta demonstracyjnego]
[CLIENT] Initializing mini-TOR connection...
[CLIENT] Proxy Node: proxy-node.local:8080
[CLIENT] Target: target-server.com:80
[CLIENT] MiniTorSocket created
[CLIENT] Connecting to target through proxy node...
[CLIENT] Tunnel established to target-server.com:80
[CLIENT] Connected to target-server.com through proxy-node.local
socket wrapped
[CLIENT] HTTP Request: GET / HTTP/1.1
Host: target-server.com
User-Agent: mini-TOR-Client
Connection: close


[CLIENT] Sending HTTP request...
[CLIENT] Waiting for response...
[CLIENT] ... received 168 bytes
[CLIENT] Connection closed by remote host.

[CLIENT] Final response received from target server:
------------------------------
HTTP/1.1 200 OK
Content-Type: text/plain
Content-Length: 54
Server: mini-TOR-Mock-Server
Connection: close

Hello from the Target Server! Your identity is hidden.
------------------------------
[CLIENT] Closing connection
\end{lstlisting}

\subsection{Test funkcjonalny: Anonimowość}
Weryfikacja, czy serwer docelowy widzi adres IP węzła (\texttt{proxy-node}), a nie klienta.

\begin{lstlisting}[caption=Logi z proxy-nodea]
proxy-node.local   | [NODE] [+] New connection from ('172.20.0.4', 41308)
proxy-node.local   | [NODE] [+] Handler thread started for ('172.20.0.4', 41308)
\end{lstlisting}

Widać, że IP klienta to 172.20.0.4

\begin{lstlisting}[caption=Logi serwera docelowego]
target-server.com  | [TARGET] ALERT: Incoming connection established!
target-server.com  | [TARGET] SOURCE IP (REMOTE_ADDR): 172.20.0.3
target-server.com  | [TARGET] SOURCE PORT: 41920
target-server.com  | [TARGET] Received Data Stream:
target-server.com  | ----------------------------------------
target-server.com  | GET / HTTP/1.1
target-server.com  | Host: target-server.com
target-server.com  | User-Agent: mini-TOR-Client
target-server.com  | Connection: close
target-server.com  | ----------------------------------------
target-server.com  | [TARGET] Response sent successfully.
target-server.com  | [TARGET] Connection closed.
target-server.com  | --------------------------------------------------
\end{lstlisting}

Serwer docelowy widzi adres 172.20.0.3 - adres proxy-nodea

\subsection{Test funkcjonalny: Traffic Shaping}
Demonstracja segmentacji danych. Klient otrzymuje odpowiedź w "paczkach" z~widocznymi opóźnieniami.

Zmodyfikowaliśmy metodę recv klasy DoubleSocket w celu podejrzenia czasu opóźnień

\begin{lstlisting}[caption=Modyfikacja metody]
    def recv(self, n):
        while True:
            try:
                time_start = time.time()
                encrypted = self.sock.recv(n)
                delay = time.time() - time_start
                print(f"Delay for recv: {delay:.4f} seconds")
                if not encrypted:
                    return encrypted
                self._incoming.write(encrypted)
                return self._obj.read()
            except ssl.SSLWantReadError:
                pass

\end{lstlisting}

\begin{lstlisting}[caption=Logi klienta demonstracyjnego (Segmentacja)]
[CLIENT] Sending HTTP request...
[CLIENT] Waiting for response...
Delay for recv: 0.1839 seconds
Delay for recv: 0.0810 seconds
Delay for recv: 0.0000 seconds
Delay for recv: 0.0000 seconds
Delay for recv: 0.1452 seconds
Delay for recv: 0.1749 seconds
Delay for recv: 0.0000 seconds
Delay for recv: 0.0729 seconds
Delay for recv: 0.0000 seconds
Delay for recv: 0.0000 seconds
Delay for recv: 0.0000 seconds
[CLIENT] ... received 335 bytes
Delay for recv: 0.0018 seconds
[CLIENT] Connection closed by remote host.
\end{lstlisting}

Widać, że czas opóźnienia jest pseudolosowy

\section{Podsumowanie}
Zrealizowany projekt spełnia założenia dokumentacji wstępnej.
\begin{itemize}
    \item \textbf{Anonimowość:} Została potwierdzona testami; serwer docelowy nie ma wiedzy o oryginalnym adresie IP klienta.
    \item \textbf{Bezpieczeństwo:} Komunikacja na odcinku Klient-Węzeł jest zabezpieczona protokołem TLS, co chroni przed podsłuchem w sieci lokalnej klienta.
    \item \textbf{Analiza ruchu:} Węzeł skutecznie wprowadza losowe opóźnienia i fragmentację, co utrudnia analizę czasową przepływu danych.
\end{itemize}

Rozwiązanie stanowi funkcjonalny prototyp biblioteki anonimizującej ("Mini-TOR").

\end{document}